%% sections/background.tex
\section{Background: Konkani as a Low-Resource Language}
\label{sec:background}

\subsection{Socio-cultural and Linguistic Profile}

Konkani (\textit{Koṅkaṇī}) is an Indo-Aryan language belonging to the Indo-Iranian branch of the Indo-European family. It is the sole official language of the state of Goa and is also spoken in the coastal districts of Karnataka, Kerala, and Maharashtra. The 2011 Census of India recorded approximately 2.3 million mother-tongue speakers, though estimates including second-language speakers place the community considerably larger.

The language occupies a paradoxical position in the Indian linguistic landscape. On the one hand, it holds the status of a Scheduled Language under the Eighth Schedule of the Indian Constitution—one of only 22 languages to do so—and was recognised as the official language of Goa in 1987. On the other hand, it has historically been overshadowed by Marathi (the dominant language of the surrounding region) and Portuguese (the colonial language of Goa for over 450 years), which has led to a fragmented and under-standardised written tradition.

Konkani is characterised by considerable dialectal variation. Major dialects include Antruzi (the basis for the standardised literary register), Bardezi, Saxtti, and Malvani, among others. Written Konkani appears in multiple scripts: Devanagari (the official and most widely used script in Goa), Kannada script (in Karnataka), Malayalam script (in Kerala), and the Roman script (historically used by the Goan Catholic community). This orthographic diversity complicates the collection and processing of written data at scale.

Linguistically, Konkani is a morphologically rich language with a relatively free word order (predominantly SOV). Its nominal morphology encodes gender (masculine, feminine, neuter), number, and case. Verbal morphology is complex, encoding tense, aspect, mood, and agreement with subject and, in some constructions, object. These properties pose challenges for both tokenisation and the learning of syntactic and semantic structure from limited data.

\subsection{Low-Resource Nature and NLP Challenges}

Despite its constitutional recognition, Konkani is severely under-resourced by every standard metric used in the NLP community:

\textbf{Corpora.} Large-scale monolingual corpora do not exist for Konkani. Unlike Hindi or Bengali, which are represented in Common Crawl, Wikipedia, and numerous parallel corpora, Konkani's web footprint is limited. Wikipedia in Konkani (\texttt{gom.wikipedia.org}) contains roughly 3,800 articles, several orders of magnitude smaller than high-resource Indic languages.

\textbf{Annotated data.} There are no publicly available treebanks, named entity corpora, or semantic banks for Konkani. Universal Dependencies does not include a Konkani treebank. No AMR dataset exists prior to this work.

\textbf{Pre-trained models.} Konkani is absent from most large pre-trained multilingual models. The mBART-50 tokenizer, which covers 50 languages, does not include Konkani; the closest available language code is Hindi (\texttt{hi\_IN}). IndicBERT \cite{kakwani2020indicnlpsuite} and IndicTrans2 \cite{gala2023indictrans2} include Konkani among the 22 scheduled languages, making them among the few exceptions.

\textbf{Translation resources.} The AI4Bharat BPCC dataset \cite{gala2023indictrans2} provides parallel Konkani–English sentence pairs mined from government documents, news, and other public sources. This is the primary parallel resource available for Konkani and is used in this work for both annotation and pretraining.

\subsection{Existing Resources and the Need for Semantic Resources}

The landscape of Konkani NLP resources can be summarised as follows. Machine translation quality has improved with the release of IndicTrans2, which supports Konkani as a target language. Basic tasks such as language identification and script normalisation are partially addressed. However, higher-level semantic processing—dependency parsing, semantic role labelling, named entity recognition, and above all semantic representation—remains entirely unstudied for Konkani.

Abstract Meaning Representation is particularly valuable in the low-resource setting because it provides a language-independent semantic scaffold. An AMR graph for a Konkani sentence refers to the same PropBank-style predicate senses and semantic relations as English AMR, making cross-lingual transfer and evaluation possible. Moreover, AMR enables the extraction of structured knowledge from unstructured Konkani text, which could unlock downstream applications in information retrieval, dialogue systems, and question answering for the Konkani-speaking community.

Creating the first AMR dataset for Konkani is thus both a practical and a scientific contribution: it directly advances the resource landscape for an under-served community while providing a testbed for research on low-resource semantic parsing more broadly.
